\chapter{Reference Guide}
\label{toc:reference}

This chapter provides a detailed description of the side conditions generated by
the \Region plug-in, via its command-line or via its programmatic interface,
together with the associated optimizations and relevant command-line options. It
is organized by source of side-conditions: Section~\ref{toc:code} for
expression, l-values and \C instructions and Section~\ref{toc:functions} for
function pre-conditions. Special cases modify the generated side-conditions as
follows:
\begin{description}
  \item[skipped:] the annotation is not generated since it is statically verified;
  \item[invalid:] the annotation is assigned a False-if-reachable status;
\end{description}

Annotations generated through the command-line are tagged with \verb"region"
name. Invalid annotations are tagged with an additional \verb"invalid" property
name. Each category of annotations has also dedicated property names.

%% --------------------------------------------------------------------------
\section{Side Conditions}
\label{toc:code}
%% --------------------------------------------------------------------------

%% --------------------------------------------------------------------------
\subsection{Bounds}

\begin{lstlisting}
//@ check region: 0 <= k < N;
return a[k];
\end{lstlisting}

\noindent
Generated when accessing an array element, unless option \verb+-unsafe-arrays+
is set.

%% --------------------------------------------------------------------------
\subsection{Valid Region}

\begin{lstlisting}
//@ check region: path: \valid_region(p+k, sizeof(int));
return *(p+k);
\end{lstlisting}

\noindent
Generated when accessing a region with pointer arithmetic. Also generated when
accessing array elements when option \verb+-unsafe-arrays+ is set.

%% --------------------------------------------------------------------------
\subsection{Separated Objects}

\begin{lstlisting}
//@ check region: f: A: B: \separated(a,b);
r = f(a,b);
\end{lstlisting}

\noindent
Generated when two distinct objects \verb+A+ and \verb+B+ accessible from function
\verb+f+ might be aliased at a call site.

\subsection{Others}
\TODO{To be complete}
%% --------------------------------------------------------------------------
